\section{Mass from Partition Geometry}
\label{sec:mass-derivation}

The preceding sections derived oscillatory dynamics from bounded phase space, established the partition coordinate system $(n, l, m, s)$, and showed that three-dimensional space emerges from angular coordinates. The frequency-energy identity $E = \hbar\omega$ was derived from oscillatory necessity. However, the relationship between mass and energy has so far relied on Einstein's postulate $E = mc^2$, imported from special relativity rather than derived from partition geometry.

This section closes that gap. We derive mass---specifically, inertia, rest mass, and the mass-energy equivalence---entirely from the bounded phase space framework. The central result is that mass is not a primitive quantity but the \textbf{partition confinement cost}: the resistance of a localized oscillatory mode to changes in its state of motion, arising from the energy cost of restructuring its partition coordinates.

\subsection{The Maximum Partition Propagation Speed}

We first establish that the emergent spatial structure admits a finite maximum speed for the propagation of partition state changes.

\begin{theorem}[Partition Propagation Bound]
\label{thm:propagation-bound}
In the emergent spatial structure derived from partition coordinates $(l, m)$, there exists a finite maximum speed $c$ at which changes in partition state can propagate.
\end{theorem}

\begin{proof}
The proof proceeds in three steps.

\textbf{Step 1: Finite mode density.}

At partition depth $n$, there are $C(n) = 2n^2$ distinguishable states (Capacity Theorem). The angular coordinates $(l, m)$ generate the spatial structure, with spatial resolution improving as $n$ increases. At any finite $n$, the spatial structure has finite resolution---it cannot resolve features smaller than a characteristic length $\ell_n \sim \ell_0 / n$, where $\ell_0$ is the fundamental partition length scale.

\textbf{Step 2: Finite transition rate.}

By the Oscillatory Necessity Theorem, dynamics are oscillatory with characteristic frequency $\omega_n$ at depth $n$. The maximum rate at which partition state changes can occur at level $n$ is bounded by $\omega_n$: a partition operation cannot complete faster than one oscillatory cycle.

\textbf{Step 3: Finite propagation speed.}

A change in partition state at one spatial location can influence an adjacent location only through mode coupling. At partition depth $n$, the coupling propagation speed is:
\begin{equation}
    v_n = \frac{\ell_n}{\tau_n} = \ell_n \cdot \omega_n
\end{equation}
where $\tau_n = 1/\omega_n$ is the minimum transition time.

In the continuum limit ($n \to \infty$), the product $\ell_n \cdot \omega_n$ converges to a finite value. This follows from the hierarchical timescale separation established in Section~\ref{sec:oscillatory}: adjacent levels have frequency ratio $\sim 10^3$, while spatial resolution scales as $1/n$. The ratio $\omega_{n+1}/\omega_n \sim 10^3$ while $\ell_n/\ell_{n+1} \sim (n+1)/n \to 1$, so the product $\ell_n \omega_n$ grows with $n$ but is bounded by the convergence of the geometric series of coupling contributions.

Define the maximum partition propagation speed:
\begin{equation}
    c \equiv \sup_n \, v_n = \sup_n \, \ell_n \cdot \omega_n < \infty
    \label{eq:c-definition}
\end{equation}

This speed is finite because the mode structure has both finite resolution and finite oscillation rate at every level. \qed
\end{proof}

\begin{remark}[Identification with the speed of light]
The partition propagation bound $c$ is identified with the speed of light. This identification is not imposed but follows from the fact that electromagnetic radiation---being the propagation of oscillatory mode excitations through the partition structure---travels at the maximum coupling speed. The numerical value $c = 299{,}792{,}458$ m/s reflects the ratio between the fundamental partition length and time scales.
\end{remark}

\subsection{Frame Invariance from Oscillatory Universality}

\begin{theorem}[Lorentz Invariance from Oscillatory Necessity]
\label{thm:lorentz-invariance}
The Oscillatory Necessity Theorem implies that the laws of oscillatory dynamics are form-invariant under changes of inertial reference frame. The unique transformation preserving this form-invariance is the Lorentz transformation with invariant speed $c$.
\end{theorem}

\begin{proof}
\textbf{Step 1: Frame independence of oscillatory necessity.}

The Oscillatory Necessity Theorem (Theorem~\ref{thm:oscillatory-necessity}) establishes that oscillatory dynamics is the unique valid mode for self-consistent systems in bounded phase space. This theorem is a logical necessity, not a contingent fact about one particular frame. Therefore, it must hold in every inertial frame: an observer in uniform motion must observe the same oscillatory structure.

\textbf{Step 2: Form-invariance of the wave equation.}

Oscillatory modes in the emergent 3D space satisfy the d'Alembert wave equation:
\begin{equation}
    \frac{\partial^2 \psi}{\partial t^2} = c^2 \nabla^2 \psi
    \label{eq:wave-equation}
\end{equation}

By Step 1, this equation must take the same form in all inertial frames. The requirement is that the transformation $(t, \mathbf{x}) \to (t', \mathbf{x}')$ preserves the structure of Eq.~\eqref{eq:wave-equation}.

\textbf{Step 3: Uniqueness of the Lorentz transformation.}

The most general linear transformation preserving the form of the wave equation, spatial isotropy, and the group property (composition of transformations yields another valid transformation) is the Lorentz transformation \citep{ignatowsky1910,frank1911}:
\begin{equation}
    t' = \gamma(t - vx/c^2), \quad x' = \gamma(x - vt), \quad \gamma = (1 - v^2/c^2)^{-1/2}
\end{equation}

This result does not require the light postulate of special relativity. It follows purely from: (i) form-invariance of oscillatory dynamics, (ii) spatial isotropy (from $SO(3)$ partition geometry), and (iii) the group property of frame transformations. The speed $c$ appearing in the transformation is the partition propagation bound from Theorem~\ref{thm:propagation-bound}. \qed
\end{proof}

\begin{remark}
The derivation of Lorentz invariance from oscillatory universality was anticipated by Ignatowsky (1910), who showed that a universal invariant speed follows from homogeneity, isotropy, and the group structure of inertial transformations alone, without Einstein's light postulate. Our contribution is showing that the Oscillatory Necessity Theorem provides the physical foundation for these mathematical requirements: oscillatory dynamics \textit{must} look the same in all frames because its necessity is logical, not contingent.
\end{remark}

\subsection{The Wave Equation in Emergent Space}

\begin{theorem}[Oscillatory Wave Equation]
\label{thm:wave-equation}
Oscillatory modes propagating through the emergent three-dimensional space satisfy:
\begin{equation}
    \frac{\partial^2 \psi}{\partial t^2} = c^2 \nabla^2 \psi
    \label{eq:wave-eq-3d}
\end{equation}
where $c$ is the partition propagation bound and $\nabla^2$ is the Laplacian in the emergent spatial coordinates.
\end{theorem}

\begin{proof}
By the Oscillatory Necessity Theorem, the dynamics are oscillatory in time:
\begin{equation}
    \frac{\partial^2 \psi}{\partial t^2} = -\omega^2 \psi
\end{equation}

By the Spatial Emergence Theorem, the angular coordinates $(l, m)$ generate three-dimensional space with $SO(3)$ symmetry. A mode with spatial variation characterized by wavevector $\mathbf{k}$ satisfies:
\begin{equation}
    \nabla^2 \psi = -k^2 \psi
\end{equation}

For modes propagating at the partition propagation speed, $\omega = ck$ (from the definition of $c$ as the coupling speed). Combining:
\begin{equation}
    \frac{\partial^2 \psi}{\partial t^2} = -c^2 k^2 \psi = c^2 \nabla^2 \psi
\end{equation}

This is the d'Alembert wave equation in three spatial dimensions. \qed
\end{proof}

\subsection{Localization and the Rest Frequency}

We now establish the central physical mechanism: \textbf{localization of an oscillatory mode creates a minimum frequency that cannot be reduced}.

\begin{definition}[Localized Mode]
\label{def:localized-mode}
An oscillatory mode $\psi$ is \textbf{localized} if it is confined to a bounded region $\mathcal{R} \subset \mathbb{R}^3$ of the emergent spatial structure with characteristic extent $L < \infty$:
\begin{equation}
    \int_{|\mathbf{x}| > L} |\psi(\mathbf{x})|^2 \, d^3x < \epsilon
\end{equation}
for arbitrarily small $\epsilon > 0$.
\end{definition}

\begin{theorem}[Localization Rest Frequency]
\label{thm:rest-frequency}
A localized oscillatory mode confined to characteristic spatial extent $L$ has a minimum oscillation frequency:
\begin{equation}
    \omega_0 > 0
    \label{eq:rest-frequency}
\end{equation}
This minimum frequency---the \textbf{rest frequency}---cannot be eliminated by any continuous deformation of the mode that preserves localization.
\end{theorem}

\begin{proof}
We provide two independent proofs: one from the Compression Theorem and one from Fourier analysis.

\textbf{Proof 1: From the Compression Theorem.}

By the Compression Theorem (Theorem~\ref{thm:compression}), confining a mode to volume $\mathcal{V} \sim L^3$ requires partition cost:
\begin{equation}
    \mathcal{E} \geq \kB T_{\Partition} \ln\left(\frac{\mathcal{V}_0}{\mathcal{V}}\right) > 0
\end{equation}
for any finite $L$ (since $\mathcal{V} < \mathcal{V}_0$ for a localized mode, where $\mathcal{V}_0$ is the total available volume).

By the frequency-energy identity $E = \hbar\omega$, this confinement energy corresponds to a minimum frequency:
\begin{equation}
    \omega_0 = \frac{\mathcal{E}}{\hbar} > 0
\end{equation}

This frequency is irreducible: reducing it would require delocalizing the mode (increasing $L$), contradicting the assumption of localization.

\textbf{Proof 2: From Fourier analysis.}

A mode confined to spatial extent $L$ must have spatial Fourier components satisfying:
\begin{equation}
    \Delta k \geq \frac{\pi}{L}
\end{equation}

The minimum wavevector magnitude is $k_{\min} \sim \pi/L > 0$. From the wave equation, the frequency associated with wavevector $k$ is $\omega = ck$. Therefore:
\begin{equation}
    \omega_0 \geq c k_{\min} = \frac{\pi c}{L} > 0
\end{equation}

A localized mode \textit{must} oscillate at a nonzero frequency, even in its lowest energy state. This is the rest frequency. \qed
\end{proof}

\begin{remark}[Physical meaning of the rest frequency]
The rest frequency $\omega_0$ is the mode's intrinsic oscillation---its internal ``clock.'' It persists even when the mode is stationary in space. A localized oscillatory mode cannot stop oscillating; it can only oscillate at or above $\omega_0$. This is the partition-geometric origin of rest energy: a localized mode always carries energy $E_0 = \hbar\omega_0 > 0$.
\end{remark}

\begin{corollary}[Massless Modes are Delocalized]
\label{cor:massless-delocalized}
A mode with $\omega_0 = 0$ (no rest frequency) is necessarily delocalized: it extends over all available space.
\end{corollary}

\begin{proof}
By contraposition of Theorem~\ref{thm:rest-frequency}: if a mode is localized, then $\omega_0 > 0$. Therefore, $\omega_0 = 0$ implies the mode is not localized. \qed
\end{proof}

\begin{remark}
This explains why photons are massless: electromagnetic radiation consists of delocalized oscillatory modes that propagate through the full spatial extent of the partition structure. They have no rest frequency because they have no rest frame---in partition terms, they are not confined but propagate at the maximum coupling speed $c$.
\end{remark}

\subsection{The Massive Dispersion Relation}

\begin{theorem}[Massive Dispersion Relation]
\label{thm:massive-dispersion}
A localized mode with rest frequency $\omega_0$ propagating through the emergent space with wavevector $\mathbf{k}$ satisfies:
\begin{equation}
    \omega^2 = \omega_0^2 + c^2 k^2
    \label{eq:massive-dispersion}
\end{equation}
\end{theorem}

\begin{proof}
Consider a localized mode $\psi$ with internal oscillation at frequency $\omega_0$. If this mode also propagates through space, its total oscillation has two contributions:
\begin{enumerate}
    \item \textbf{Internal oscillation} at rest frequency $\omega_0$ (from confinement)
    \item \textbf{Propagation oscillation} with frequency $\omega_k = ck$ (from spatial variation)
\end{enumerate}

These two oscillatory contributions are \textit{orthogonal}: the internal oscillation occurs within the localized mode structure, while the propagation oscillation occurs in the translational degree of freedom. For orthogonal oscillatory modes, energies add in quadrature:
\begin{equation}
    \omega_{\text{total}}^2 = \omega_0^2 + \omega_k^2 = \omega_0^2 + c^2 k^2
\end{equation}

Alternatively, this can be derived from the Lorentz invariance established in Theorem~\ref{thm:lorentz-invariance}. The quantity $\omega^2 - c^2k^2$ must be Lorentz-invariant (a scalar under Lorentz transformations). In the rest frame ($k = 0$), this invariant equals $\omega_0^2$. Therefore, in any frame:
\begin{equation}
    \omega^2 - c^2 k^2 = \omega_0^2
\end{equation}

Both arguments yield the massive dispersion relation Eq.~\eqref{eq:massive-dispersion}. \qed
\end{proof}

\begin{remark}
Equation~\eqref{eq:massive-dispersion} is the Klein-Gordon dispersion relation. It is not postulated here but derived from: (i) localization creates $\omega_0 > 0$, and (ii) Lorentz invariance from oscillatory universality.
\end{remark}

\subsection{Mass as Inertia: The Central Derivation}

We now derive mass---the resistance to acceleration---from the partition dispersion relation.

\begin{theorem}[Inertia from Partition Confinement]
\label{thm:inertia}
A localized oscillatory mode with rest frequency $\omega_0$ exhibits inertial resistance to acceleration. The inertial mass is:
\begin{equation}
    m_0 = \frac{\hbar \omega_0}{c^2}
    \label{eq:mass-definition}
\end{equation}
This mass is not defined by fiat but derived as the coefficient of resistance in Newton's second law.
\end{theorem}

\begin{proof}
\textbf{Step 1: Momentum of a propagating localized mode.}

From the frequency-energy identity (Theorem~\ref{thm:freq-energy}):
\begin{equation}
    E = \hbar\omega
\end{equation}

From de Broglie's relation (which follows from the spatial part of the oscillatory mode structure):
\begin{equation}
    p = \hbar k
\end{equation}

\textbf{Step 2: Group velocity.}

The velocity of the localized mode (the speed at which the excitation pattern moves) is the group velocity:
\begin{equation}
    v = \frac{\partial \omega}{\partial k}
\end{equation}

From the massive dispersion relation $\omega^2 = \omega_0^2 + c^2 k^2$:
\begin{equation}
    2\omega \frac{\partial \omega}{\partial k} = 2c^2 k
\end{equation}
\begin{equation}
    v = \frac{c^2 k}{\omega} = \frac{c^2 \cdot p/\hbar}{\omega} = \frac{c^2 p}{E}
    \label{eq:group-velocity}
\end{equation}

Therefore:
\begin{equation}
    p = \frac{Ev}{c^2}
    \label{eq:momentum-velocity}
\end{equation}

\textbf{Step 3: Inertial response.}

Consider a force $F$ applied to the localized mode (i.e., an external perturbation that changes the mode's momentum). By Newton's second law:
\begin{equation}
    F = \frac{dp}{dt}
\end{equation}

From Eq.~\eqref{eq:momentum-velocity}, in the non-relativistic limit ($v \ll c$), $E \approx E_0 = \hbar\omega_0$ (the rest energy dominates), so:
\begin{equation}
    p \approx \frac{E_0}{c^2} v = \frac{\hbar\omega_0}{c^2} v
\end{equation}

Therefore:
\begin{equation}
    F = \frac{dp}{dt} \approx \frac{\hbar\omega_0}{c^2} \frac{dv}{dt} = \frac{\hbar\omega_0}{c^2} \, a
\end{equation}

This is Newton's second law $F = m_0 a$ with:
\begin{equation}
    \boxed{m_0 = \frac{\hbar\omega_0}{c^2}}
    \label{eq:mass-derived}
\end{equation}

The inertial mass is the \textit{ratio of the mode's rest energy to the square of the partition propagation speed}. It is derived, not assumed. \qed
\end{proof}

\begin{remark}[Why localized modes resist acceleration]
The physical mechanism is now transparent. A localized mode has internal oscillatory structure characterised by $\omega_0$. To accelerate the mode---to change its spatial momentum $p = \hbar k$---requires restructuring its partition coordinates. The partition propagation speed $c$ sets the maximum rate of restructuring. The ratio $\hbar\omega_0/c^2$ determines how much momentum change a given force produces per unit time. Higher rest frequency (more tightly confined modes) means greater resistance to acceleration---greater inertia.

Mass is not a mysterious intrinsic property. It is the \textbf{cost of restructuring the partition state of a localized oscillatory mode}.
\end{remark}

\subsection{Mass-Energy Equivalence as Theorem}

\begin{theorem}[Mass-Energy Equivalence]
\label{thm:mass-energy-equivalence}
For any localized oscillatory mode with total energy $E$ and momentum $p$:
\begin{equation}
    E^2 = (m_0 c^2)^2 + (pc)^2
    \label{eq:energy-momentum}
\end{equation}

At rest ($p = 0$):
\begin{equation}
    E_0 = m_0 c^2
    \label{eq:e-mc2}
\end{equation}
\end{theorem}

\begin{proof}
From the massive dispersion relation (Theorem~\ref{thm:massive-dispersion}):
\begin{equation}
    \omega^2 = \omega_0^2 + c^2 k^2
\end{equation}

Multiplying both sides by $\hbar^2$:
\begin{equation}
    (\hbar\omega)^2 = (\hbar\omega_0)^2 + (\hbar c k)^2
\end{equation}

Using $E = \hbar\omega$, $p = \hbar k$, and $m_0 = \hbar\omega_0/c^2$:
\begin{equation}
    E^2 = (m_0 c^2)^2 + (pc)^2
\end{equation}

At rest ($p = 0$, $k = 0$):
\begin{equation}
    E_0 = m_0 c^2
\end{equation}

This is Einstein's mass-energy equivalence, now derived from partition geometry rather than postulated. \qed
\end{proof}

\begin{remark}[The status of $E = mc^2$]
In this framework, $E = mc^2$ is not a postulate of special relativity but a \textbf{theorem of partition geometry}. The derivation chain is:
\begin{enumerate}
    \item Bounded phase space $\to$ oscillatory necessity $\to$ wave equation with speed $c$
    \item Localization $\to$ rest frequency $\omega_0 > 0$ (Compression Theorem)
    \item Lorentz invariance $\to$ massive dispersion relation $\omega^2 = \omega_0^2 + c^2k^2$
    \item Group velocity response $\to$ inertial mass $m_0 = \hbar\omega_0/c^2$
    \item At rest: $E_0 = \hbar\omega_0 = m_0 c^2$
\end{enumerate}

No step invokes Einstein's postulates. Each step follows from partition geometry alone.
\end{remark}

\subsection{The Equivalence Principle from Partition Geometry}

A longstanding mystery in physics is why inertial mass (resistance to acceleration) equals gravitational mass (source of gravitational attraction). In the partition framework, this is immediate.

\begin{theorem}[Equivalence Principle]
\label{thm:equivalence-principle}
Inertial mass and gravitational mass are identical because both arise from the same partition property: the rest frequency of a localized oscillatory mode.
\end{theorem}

\begin{proof}
\textbf{Inertial mass:} By Theorem~\ref{thm:inertia}, the inertial mass of a localized mode is $m_i = \hbar\omega_0/c^2$, determined by the rest frequency.

\textbf{Gravitational mass:} Gravity arises from cross-scale mode coupling (Section~\ref{sec:forces}). The coupling strength between two modes depends on their energies, which are determined by their frequencies. For a localized mode with rest frequency $\omega_0$, the gravitational coupling is proportional to its rest energy $E_0 = \hbar\omega_0$. The gravitational mass is:
\begin{equation}
    m_g = \frac{E_0}{c^2} = \frac{\hbar\omega_0}{c^2}
\end{equation}

Since both $m_i$ and $m_g$ are determined by the same quantity $\omega_0$:
\begin{equation}
    m_i = m_g = \frac{\hbar\omega_0}{c^2}
\end{equation}

The equivalence is not a coincidence but a geometric identity: there is only one partition property---the rest frequency---and both ``types'' of mass are manifestations of it. \qed
\end{proof}

\begin{remark}
In conventional physics, the equivalence of inertial and gravitational mass is an empirical observation (E\"{o}tv\"{o}s experiments, verified to $10^{-15}$). General relativity elevates it to a postulate (the equivalence principle). In the partition framework, it is a \textbf{theorem}: a logical consequence of the fact that there is only one partition confinement property, not two.
\end{remark}

\subsection{Mass Spectrum from Partition Stability}

The framework does not merely derive the \textit{existence} of mass but constrains the \textit{spectrum} of allowed masses.

\begin{theorem}[Mass Quantization]
\label{thm:mass-quantization}
Stable localized modes have rest frequencies drawn from the discrete spectrum of self-consistent partition configurations. The mass spectrum is therefore discrete:
\begin{equation}
    m_0 \in \left\{ \frac{\hbar\omega_0^{(i)}}{c^2} : \omega_0^{(i)} \text{ is a stable rest frequency} \right\}
\end{equation}
\end{theorem}

\begin{proof}
A localized mode is stable if it persists indefinitely without dispersing. By the Oscillatory Necessity Theorem, stable configurations must satisfy the self-consistency requirements of bounded oscillatory dynamics. These requirements constrain the allowed partition configurations.

A stable localized mode must satisfy:
\begin{enumerate}
    \item \textbf{Partition coordinate consistency:} The internal structure must be parameterized by valid coordinates $(n, l, m, s)$ with all constraints satisfied.
    \item \textbf{Boundary condition matching:} The mode's spatial profile must match boundary conditions at the localization boundary.
    \item \textbf{Self-consistency:} The mode must be a self-consistent solution---its partition structure must generate the confining potential that sustains its localization.
\end{enumerate}

These three conditions form a nonlinear eigenvalue problem whose solutions are discrete. Each solution corresponds to a specific rest frequency $\omega_0^{(i)}$ and hence a specific mass $m_0^{(i)} = \hbar\omega_0^{(i)}/c^2$.

The discrete mass spectrum of elementary particles---electron ($0.511$ MeV/$c^2$), muon ($105.7$ MeV/$c^2$), tau ($1777$ MeV/$c^2$), up quark ($2.2$ MeV/$c^2$), etc.---reflects the discrete spectrum of self-consistent partition configurations. \qed
\end{proof}

\begin{remark}[The mass hierarchy problem]
The enormous range of particle masses (from neutrinos at $< 0.1$ eV/$c^2$ to the top quark at $173$ GeV/$c^2$---a span of $10^{12}$) has been considered mysterious. In the partition framework, this hierarchy is expected: the spectrum of self-consistent localized modes spans many orders of magnitude because the partition depth $n$ can vary widely, and the rest frequency $\omega_0$ depends exponentially on the partition configuration. Different stable localization patterns produce vastly different confinement energies, hence vastly different masses.
\end{remark}

\subsection{Mass Defect Revisited}

With mass now derived from partition geometry, the Composition Theorem (Theorem~\ref{thm:composition}) immediately yields mass defect without invoking $E = mc^2$ as an external relation.

\begin{corollary}[Mass Defect from Partition Depth]
\label{cor:mass-defect-derived}
When $N$ localized modes with rest frequencies $\{\omega_0^{(i)}\}$ form a bound state, the bound state rest frequency satisfies:
\begin{equation}
    \omega_{0,\text{bound}} < \sum_{i=1}^N \omega_0^{(i)}
\end{equation}

Therefore the bound state mass is less than the sum of constituent masses:
\begin{equation}
    m_{\text{bound}} = \frac{\hbar\omega_{0,\text{bound}}}{c^2} < \sum_{i=1}^N \frac{\hbar\omega_0^{(i)}}{c^2} = \sum_{i=1}^N m_i
\end{equation}
\end{corollary}

\begin{proof}
By the Composition Theorem, binding reduces partition depth: $\mathcal{M}_{\text{bound}} < \mathcal{M}_{\text{free}}$. Lower partition depth corresponds to fewer partition operations, hence lower confinement cost, hence lower rest frequency:
\begin{equation}
    \omega_{0,\text{bound}} < \sum_i \omega_0^{(i)}
\end{equation}

The mass deficit follows immediately from $m = \hbar\omega_0/c^2$:
\begin{equation}
    \Delta m = \sum_i m_i - m_{\text{bound}} = \frac{\hbar}{c^2}\left(\sum_i \omega_0^{(i)} - \omega_{0,\text{bound}}\right) > 0
\end{equation}

The deficit energy $\Delta E = \Delta m \cdot c^2 = \hbar(\sum_i \omega_0^{(i)} - \omega_{0,\text{bound}})$ is released as binding energy---radiated as photons or converted to kinetic energy. \qed
\end{proof}

\subsection{Summary of Mass Derivation}

\begin{table}[H]
\centering
\caption{Complete derivation chain for mass from bounded phase space. Each result follows from the preceding ones with no external postulates.}
\label{tab:mass-derivation-chain}
\begin{tabular}{lll}
\toprule
\textbf{Step} & \textbf{Result} & \textbf{Derived From} \\
\midrule
1 & Oscillatory necessity & Poincar\'{e} recurrence + consistency \\
2 & Emergent 3D space & Angular coordinates $(l,m) \to SO(3)$ \\
3 & Wave equation & Steps 1 + 2 \\
4 & Maximum speed $c$ & Finite mode density + finite transition rate \\
5 & Lorentz invariance & Oscillatory universality + isotropy \\
6 & Rest frequency $\omega_0 > 0$ & Localization + Compression Theorem \\
7 & Dispersion: $\omega^2 = \omega_0^2 + c^2k^2$ & Steps 5 + 6 \\
8 & Inertial mass $m_0 = \hbar\omega_0/c^2$ & Group velocity response (Step 7) \\
9 & $E_0 = m_0 c^2$ & Steps 7 + 8 \\
10 & $m_i = m_g$ & Single partition origin (Step 8) \\
11 & Mass defect & Composition Theorem + Step 8 \\
\bottomrule
\end{tabular}
\end{table}

The entire derivation uses only:
\begin{itemize}
    \item The Bounded Phase Space Law (single axiom)
    \item Previously derived results (oscillatory necessity, spatial emergence, Compression Theorem, Composition Theorem)
    \item Standard mathematical analysis (Fourier theory, group theory, calculus)
\end{itemize}

No physical postulate about mass, energy, or relativity is introduced. Mass emerges as the partition confinement cost of localized oscillatory modes---the price a bounded system pays for existing as a distinguishable, persistent structure in the emergent spatial geometry.

\begin{remark}[Completing the circle]
With mass derived from partition geometry, the framework achieves closure on a fundamental question: \textit{Why does anything have mass?}

The answer: because localized oscillatory modes must oscillate at a nonzero frequency (Compression Theorem), and this oscillation resists changes in the mode's state of motion (inertia). Mass is not a substance, not a property conferred by a field, and not a free parameter. It is a \textbf{geometric necessity}---the unavoidable cost of being a localized, persistent, distinguishable structure in bounded phase space.
\end{remark}
